\section{Prime Contacts}
The height has an exact arithmetic meaning at integer inputs. Each active
wave contributes one for a discovered proper divisor and zero otherwise.
This identifies the contacts without making divisor counting part of the
recurrence itself.

\begin{itemize}
\item Truncated proper-divisor counts at finite stages.
\item Divisor height and exact prime contacts in the limit.
\end{itemize}

Let $n\ge2$ be an integer. Write $\mathbf 1_{m\mid n}$ for the indicator
that equals one when $m$ divides $n$, and zero otherwise. At an integer,
$a_m(n)$ is one if $m < n$ and zero if $m\ge n$. Hence
\begin{equation*}
\begin{aligned}
h_M(n)
&=\sum_{m=2}^{M}a_m(n)W_m(n)
&&[\sin(\pi n)=0]\\
&=\sum_{m=2}^{\min(M,n-1)}\mathbf 1_{m\mid n}.
&&[\text{Integer Wave Values}]
\end{aligned}
\end{equation*}

Let $\tau(n)$ count the positive divisors of $n$, including $1$ and $n$.
This is the classical divisor function, denoted $d(n)$ in DLMF \cite{dlmfdivisors}.
Once every possible proper divisor is covered, the height is
\begin{equation*}
\begin{aligned}
h(n)&=\sum_{m=2}^{n-1}\mathbf 1_{m\mid n}
&&[\text{Local Stabilization}]\\
&=\tau(n)-2.
&&[\text{By Definition}]
\end{aligned}
\end{equation*}

A prime has exactly its two trivial divisors. Every composite integer has
at least one additional divisor. Noninteger inputs cannot create extra
contacts because the baseline is strictly positive there:
\begin{equation*}
\begin{aligned}
t\notin\mathbb Z&\ \Longrightarrow\ h(t)\ge\sin^2(\pi t) > 0,\\
n\in\mathbb Z,\ n\ge2&\ \Longrightarrow\
h(n)=0\ \Longleftrightarrow\ \tau(n)=2.
\end{aligned}
\end{equation*}

Together these statements give the exact contact characterization:
\begin{equation*}
\begin{aligned}
F(t)\in\mathbb R
&\ \Longleftrightarrow\ h(t)=0\\
&\ \Longleftrightarrow\ t\text{ is prime},\qquad t\ge2.
\quad\blacksquare\ \text{[Q.E.D.]}
\end{aligned}
\end{equation*}

For example, $h(7)=0$, whereas $h(12)=4$ because $2,3,4,6$ are its
nontrivial proper divisors. The height records overlaps between filters,
while its zero set records primality.

Selecting real-axis contacts is essential. Since the real part of $F(t)$ is
always $t$, projecting the entire curve onto the real axis would retain every
input. This construction represents accepted values in value coordinates;
it does not enumerate primes by their position in a sequence.
